\newcommand{\capa}[2]{%
\begin{titlepage}\centering\singlespacing
\includegraphics[width=4.5cm]{java/img/logo.png}\\[1.2cm]
{\large UNIVERSIDADE CAT\'OLICA DO SALVADOR\\ Evolu\c{c}\~ao e Moderniza\c{c}\~ao de Sistemas Legados -- 2026.2}\\[3cm]
{\Large\bfseries #1}\\[0.8cm]
{\large #2}\\[1.2cm]
{\large Projeto: SupermarketReceipt Refactoring Kata (Java)\\ Grupo 2: \code{ShoppingCart}, \code{ProductQuantity}, \code{Product}, \code{ProductUnit}}\\[3cm]
{\large Juvenal Rosa Alves Costa Junior\\ Matr\'icula 200033214}\vfill
{\large Salvador\\ Novembro de 2026}
\end{titlepage}}

\capa{Plano de Atualização Tecnológica e Estratégias para Evolução Futura}{Etapa 3 -- Refatoração Prática e Evolução Aplicada}

\tableofcontents\newpage

\section{Introdução e escopo}

Este documento apresenta o plano de atualização tecnológica e as estratégias de evolução futura do sistema legado \textit{SupermarketReceipt Refactoring Kata}, desenvolvido em Java. O sistema simula operações de um caixa de supermercado, incluindo o gerenciamento de produtos, quantidades, aplicação de promoções e emissão de recibos.

A divisão do trabalho da Etapa 3 atribui ao Grupo 2 as classes \code{ShoppingCart}, \code{ProductQuantity}, \code{Product} e \code{ProductUnit}. O principal objetivo desse grupo foi melhorar a consistência dos dados do domínio, eliminar a duplicação de estado no carrinho de compras e estabelecer validações para os objetos que representam produtos e quantidades.

Este plano considera as tecnologias já utilizadas pelo projeto e apresenta possibilidades de evolução que podem melhorar sua manutenção, precisão numérica, qualidade e capacidade de expansão. Também são identificados os pontos de integração com o Grupo 1, responsável pelas ofertas, e com o Grupo 3, responsável pelo recibo e pela impressão.

É importante distinguir as mudanças efetivamente implementadas das propostas futuras. As alterações realizadas pelo Grupo 2 estão registradas na seção 6, enquanto as demais iniciativas são apresentadas como oportunidades de evolução, sem pressupor que já tenham sido executadas.

\section{Situação tecnológica atual}

O projeto utiliza uma estrutura baseada em Java e Maven, com testes automatizados e organização das classes de domínio em um único pacote.

\begin{table}[H]\centering\small
\begin{tabularx}{\textwidth}{>{\raggedright\arraybackslash}p{3.6cm}Y}
\toprule
\textbf{Aspecto} & \textbf{Situação observada no projeto} \\
\midrule
Linguagem & Java 17, com utilização de recursos como \textit{records} e \textit{enums}. \\
\addlinespace
Build e dependências & Maven, com gerenciamento de dependências e execução automatizada dos testes. \\
\addlinespace
Testes & JUnit Jupiter 5.11.4 e ApprovalTests. Foram acrescentados testes unitários para as validações das classes do Grupo 2. \\
\addlinespace
Integração contínua & O projeto possui configuração de GitHub Actions para execução dos testes automatizados. \\
\addlinespace
Análise estática & Não foi implementada uma ferramenta específica de análise estática durante o trabalho do Grupo 2. \\
\addlinespace
Documentação & O projeto possui documentação de apoio, mas não foi realizada uma padronização completa de Javadoc nas classes do grupo. \\
\addlinespace
Modelo de dados & Produtos e quantidades são representados por \code{Product}, \code{ProductQuantity} e \code{ProductUnit}. O tipo \code{double} continua sendo utilizado para quantidades. \\
\addlinespace
Gerenciamento do carrinho & A lista de itens é a fonte única de dados persistentes. O mapa de quantidades agregadas é construído sob demanda. \\
\addlinespace
Estrutura de pacotes & As classes permanecem no pacote \code{dojo.supermarket.model}. \\
\bottomrule
\end{tabularx}
\caption{Situação tecnológica do projeto considerando o escopo do Grupo 2.}
\end{table}

A versão Java 17 foi preservada durante a refatoração. Essa decisão evita introduzir uma atualização de ambiente que não era necessária para resolver os problemas encontrados nas classes do grupo e reduz o risco de incompatibilidade com a configuração existente de compilação e integração contínua.

O uso de \code{double} também foi mantido. Embora esse tipo possa apresentar limitações de precisão em operações decimais, sua substituição afetaria interfaces utilizadas pelos demais componentes, exigindo coordenação entre os grupos.

\section{Plano de atualização tecnológica}

As propostas de atualização foram organizadas em fases, considerando o impacto de cada mudança e sua dependência em relação aos demais componentes.

\begin{table}[H]\centering\small
\begin{tabularx}{\textwidth}{>{\raggedright\arraybackslash}p{1.4cm}>{\raggedright\arraybackslash}p{4.4cm}Y}
\toprule
\textbf{Fase} & \textbf{Ação} & \textbf{Resultado esperado e critério de aceite} \\
\midrule
0 & Consolidar os testes de validação de produtos, quantidades e unidades. & Objetos inválidos não são construídos e os comportamentos válidos continuam funcionando. \\
\addlinespace
1 & Avaliar a adoção de ferramentas de análise estática, como Checkstyle ou PMD. & Identificação de problemas de estilo, complexidade e manutenção, com critérios de qualidade definidos para o projeto. \\
\addlinespace
2 & Avaliar a adoção do Maven Wrapper. & Padronização da execução do Maven entre os integrantes, reduzindo dependências de instalações locais. \\
\addlinespace
3 & Avaliar a substituição de \code{double} por \code{BigDecimal} nas operações que exigem precisão decimal. & Representação decimal mais adequada, com testes específicos para cálculos e arredondamentos. \\
\addlinespace
4 & Avaliar a evolução do modelo de quantidades e unidades. & Regras de representação e validação mais explícitas, mantendo a compatibilidade com os componentes consumidores. \\
\addlinespace
5 & Ampliar a documentação técnica das classes de domínio. & Regras de validação, responsabilidades e limitações descritas por meio de Javadoc. \\
\addlinespace
6 & Avaliar a reorganização dos pacotes do sistema. & Separação mais clara entre modelo, ofertas, fechamento da compra e recibo, sem dependências circulares. \\
\bottomrule
\end{tabularx}
\caption{Fases propostas para a evolução tecnológica do Grupo 2.}
\end{table}

As fases 0 a 2 representam melhorias de baixo impacto em relação à estrutura atual. As fases 3 e 4 exigem maior atenção, pois podem afetar contratos utilizados pelos outros grupos. As fases 5 e 6 podem ser realizadas gradualmente, conforme a necessidade de manutenção e expansão do projeto.

\section{Estratégias para evolução futura}

\subsection{Representação decimal com maior precisão}

Atualmente, o projeto utiliza \code{double} para representar quantidades e valores relacionados aos cálculos de compra. Embora seja adequado para diversas operações numéricas, esse tipo utiliza representação binária de ponto flutuante e pode apresentar diferenças de precisão em determinadas operações decimais.

Uma evolução possível é utilizar \code{BigDecimal} nos contextos que exigem precisão decimal, especialmente nos valores monetários.

Essa mudança deve ser planejada em conjunto com os Grupos 1 e 3, pois as quantidades e os valores utilizados pelo carrinho também são consumidos pelas ofertas e pelo processamento dos recibos.

A adoção de \code{BigDecimal} exigiria revisar operações aritméticas, regras de escala, arredondamento, comparações e testes. Portanto, não deve ser realizada como uma simples substituição de tipo.

\subsection{Evolução do modelo de quantidades}

A classe \code{ProductQuantity} atualmente representa um produto e sua quantidade utilizando \code{double}. A validação implementada impede quantidades negativas, valores não finitos e quantidades fracionárias para unidades discretas.

Uma evolução futura poderia tornar a representação de quantidades mais especializada, utilizando tipos específicos para quantidades inteiras e fracionárias ou um objeto de valor responsável por essas regras.

Essa abordagem poderia tornar ainda mais explícitas as diferenças entre produtos vendidos por unidade e produtos vendidos por peso.

Entretanto, a adoção dessa solução precisa considerar a complexidade adicional e o impacto nas interfaces de ofertas e recibos. A estrutura atual atende ao escopo do projeto sem exigir uma hierarquia de classes ou abstrações adicionais.

\subsection{Evolução da agregação do carrinho}

A implementação atual utiliza a lista de \code{ProductQuantity} como fonte única de dados e calcula as quantidades agregadas quando \code{productQuantities()} é chamado.

Essa abordagem reduz o risco de inconsistência entre estruturas, pois não existe mais a necessidade de atualizar uma lista e um mapa simultaneamente.

Por outro lado, cada consulta percorre os itens armazenados no carrinho. Em cenários com um volume significativamente maior de operações, pode ser necessário avaliar alternativas de otimização.

Uma possibilidade seria utilizar uma estrutura de agregação com atualização centralizada, desde que exista uma estratégia que impeça divergências entre os dados individuais e os totais calculados.

A adoção dessa alternativa deve ser baseada em medições de desempenho, evitando introduzir complexidade prematuramente.

\subsection{Regras de domínio mais extensíveis}

As validações implementadas em \code{Product}, \code{ProductQuantity} e \code{ProductUnit} representam um primeiro passo para estabelecer regras de domínio mais explícitas.

Uma evolução futura poderia ampliar essas regras para contemplar características adicionais dos produtos, como restrições específicas de comercialização, precisão permitida para quantidades fracionárias e limites operacionais.

Essas regras devem permanecer próximas aos conceitos que representam, evitando que diferentes componentes implementem verificações duplicadas.

A utilização de objetos de valor também pode ser considerada caso o domínio se torne mais complexo.

\subsection{Organização modular do sistema}

As classes do Grupo 2 permanecem no pacote \code{dojo.supermarket.model}, assim como outras classes do projeto.

Para uma evolução futura, pode ser interessante organizar o sistema em pacotes com responsabilidades mais específicas, por exemplo:

\begin{itemize}

\item \code{model}: produtos, unidades e quantidades;

\item \code{offer}: estratégias e regras de promoções;

\item \code{checkout}: gerenciamento do carrinho e fechamento da compra;

\item \code{receipt}: representação e impressão dos recibos.

\end{itemize}

Essa separação pode melhorar a localização das responsabilidades e facilitar a manutenção. Entretanto, a movimentação das classes deve considerar suas dependências atuais, evitando dependências circulares e alterações desnecessárias nos contratos públicos.

Para o tamanho atual do projeto, recomenda-se preservar uma arquitetura de monólito modular, sem introduzir serviços independentes ou infraestrutura distribuída sem uma necessidade concreta.

\subsection{Ampliação da qualidade automatizada}

A suíte de testes foi ampliada para contemplar situações válidas e inválidas nas classes do Grupo 2.

Como evolução, recomenda-se complementar os testes unitários com testes baseados em propriedades, especialmente para operações de agregação de quantidades e validações de unidades.

Também pode ser útil incorporar ferramentas de análise estática e relatórios de cobertura ao processo de integração contínua.

Essas ferramentas permitiriam acompanhar a evolução da qualidade do código e identificar regressões durante futuras alterações.

\section{Cronograma, riscos e viabilidade}

As propostas de evolução podem ser implementadas gradualmente, priorizando mudanças que tragam benefícios sem comprometer a compatibilidade do sistema.

A adoção de ferramentas de qualidade e a ampliação da documentação possuem baixo impacto estrutural. Já a alteração do tipo numérico e a reorganização dos pacotes exigem planejamento e validação conjunta.

\begin{table}[H]\centering\small
\begin{tabularx}{\textwidth}{>{\raggedright\arraybackslash}p{4.3cm}>{\raggedright\arraybackslash}p{1.8cm}Y}
\toprule
\textbf{Risco} & \textbf{Impacto} & \textbf{Mitigação} \\
\midrule
Alteração do tipo numérico & Alto & Planejar a mudança com os demais grupos, revisar as operações e criar testes de equivalência. \\
\addlinespace
Aumento da complexidade do carrinho & Médio & Manter a implementação atual enquanto não houver evidências de problemas de desempenho. \\
\addlinespace
Incompatibilidade entre componentes & Alto & Preservar os contratos existentes e coordenar alterações de interfaces entre os grupos. \\
\addlinespace
Introdução de ferramentas de análise estática & Baixo & Registrar os problemas existentes antes de estabelecer regras obrigatórias no CI. \\
\addlinespace
Reorganização dos pacotes & Médio & Analisar as dependências entre classes antes de realizar movimentações. \\
\addlinespace
Ampliação das regras de domínio & Médio & Implementar novas validações acompanhadas de testes específicos e documentação. \\
\bottomrule
\end{tabularx}
\caption{Riscos e estratégias de mitigação para a evolução futura.}
\end{table}

A viabilidade das propostas depende principalmente da necessidade de evolução do sistema e do esforço necessário para preservar a compatibilidade. Para o contexto atual, as alterações realizadas pelo Grupo 2 são suficientes para melhorar a consistência do modelo sem exigir mudanças na plataforma tecnológica.

\section{Andamento do plano}

A tabela a seguir diferencia as ações efetivamente realizadas das propostas que permanecem como possibilidades de evolução.

\begin{table}[H]\centering\small
\begin{tabularx}{\textwidth}{>{\raggedright\arraybackslash}p{1.2cm}>{\raggedright\arraybackslash}p{3.2cm}Y}
\toprule
\textbf{Fase} & \textbf{Situação} & \textbf{Registro} \\
\midrule
0 & Concluída & Foram implementadas validações em \code{Product} e \code{ProductQuantity}, além de testes específicos para os comportamentos válidos e inválidos. \\
\addlinespace
1 & Pendente & Não foi incorporada ferramenta específica de análise estática ao projeto durante o escopo do Grupo 2. \\
\addlinespace
2 & Pendente & O Maven Wrapper não foi adotado como parte das alterações realizadas pelo grupo. \\
\addlinespace
3 & Pendente & O tipo \code{double} foi preservado para manter a compatibilidade com os demais componentes. \\
\addlinespace
4 & Parcialmente atendida & Foram encapsuladas as regras de aceitação de quantidades em \code{ProductUnit}, sem introduzir novos tipos especializados. \\
\addlinespace
5 & Pendente & Não foi realizada uma padronização completa de Javadoc nas classes do grupo. \\
\addlinespace
6 & Pendente & As classes permaneceram no pacote original \code{dojo.supermarket.model}. \\
\bottomrule
\end{tabularx}
\caption{Situação das propostas de evolução tecnológica do Grupo 2.}
\end{table}

Além das ações registradas, a implementação de \code{ShoppingCart} foi simplificada por meio da remoção do mapa persistente de quantidades. O método \code{productQuantities()} passou a calcular o resultado sob demanda, mantendo a interface utilizada pelos demais componentes.

A suíte completa de testes do projeto foi executada por meio do comando:

\begin{lstlisting}
mvn -f java/pom.xml clean test
\end{lstlisting}

A execução foi concluída sem erros, indicando que as alterações realizadas preservaram os comportamentos cobertos pelos testes existentes e pelos novos testes adicionados.

\section{Indicadores de sucesso}

Os indicadores abaixo permitem avaliar as melhorias implementadas e orientar as próximas etapas de evolução.

\begin{itemize}

\item \textbf{Eliminação da duplicação de estado:} o carrinho passou a manter somente a lista de itens como estrutura persistente, enquanto o mapa de quantidades é construído sob demanda (\textbf{atingido}).

\item \textbf{Validação dos dados de produto:} nomes inválidos e unidades nulas passaram a ser rejeitados durante a construção de \code{Product} (\textbf{atingido}).

\item \textbf{Validação das quantidades:} produtos nulos, quantidades negativas, valores não finitos e quantidades fracionárias para unidades discretas passaram a ser rejeitados (\textbf{atingido}).

\item \textbf{Encapsulamento das regras de unidade:} a compatibilidade entre quantidade e unidade passou a ser expressa por \code{ProductUnit.acceptsQuantity()} (\textbf{atingido}).

\item \textbf{Preservação dos comportamentos existentes:} a execução da suíte completa de testes foi concluída sem erros (\textbf{atingido nos testes executados}).

\item \textbf{Evolução da precisão decimal:} substituir \code{double} por \code{BigDecimal} permanece como proposta futura (\textbf{pendente}).

\item \textbf{Melhoria da automação de qualidade:} a adoção de análise estática e de novas ferramentas de build permanece como proposta futura (\textbf{pendente}).

\end{itemize}

\begin{thebibliography}{9}\small

\bibitem{feathers2013} FEATHERS, Michael. \textbf{Trabalho eficaz com código legado}. Porto Alegre: Bookman, 2013.

\bibitem{gamma2000} GAMMA, Erich et al. \textbf{Padrões de projeto}: soluções reutilizáveis de software orientado a objetos. Porto Alegre: Bookman, 2000.

\bibitem{fowler2020} FOWLER, Martin. \textbf{Refatoração}: aperfeiçoando o design de código existente. 2. ed. São Paulo: Novatec, 2020.

\bibitem{sommerville2019} SOMMERVILLE, Ian. \textbf{Engenharia de software}. 10. ed. São Paulo: Pearson Education do Brasil, 2019.

\bibitem{martin2011} MARTIN, Robert C. \textbf{Código limpo}: habilidades práticas do Agile software. Rio de Janeiro: Alta Books, 2011.

\end{thebibliography}